\documentclass[10pt,twocolumn]{article}

\usepackage[margin=0.72in]{geometry}
\usepackage{amsmath}
\usepackage{amssymb}
\usepackage{booktabs}
\usepackage{enumitem}
\usepackage{graphicx}
\usepackage{microtype}
\usepackage{xcolor}
\usepackage[round,authoryear]{natbib}
\usepackage[colorlinks=true,allcolors=blue]{hyperref}

\newcommand{\method}{\textsc{LocalReplay}}
\newcommand{\result}[1]{\textcolor{magenta}{\texttt{[[#1]]}}}
\newcommand{\js}{\operatorname{JS}}
\newcommand{\risk}{\operatorname{Risk}}
\newcommand{\coverage}{\operatorname{Cov}}

\title{Commutation Is Not Certification}
\author{Anonymous workshop submission}
\date{}

\begin{document}
\maketitle

\begin{abstract}
Cross-model key--value (KV) cache translation can avoid recomputing a long
context when an inference workflow switches language models.  Existing work
primarily characterizes whether a translator and model pair work on average;
it does not by itself tell a serving system whether one particular translated
cache should be trusted.  We formulate \emph{selective cross-model cache
reuse}: using only source-model state, translated state, and bounded
target-model computation, decide whether to accept a handoff or fall back to
native target prefill.  We study a diagnostic that compares two local routes:
translate an old prefix before target-side suffix replay, or process that
suffix with the source before translation.  Their discrepancy is a
translation--transition commutator.  We prove that zero commutator is not a
certificate of native-state fidelity: it measures how translation error
evolves under the probed transition and can miss harmful error shared by both
routes.  We evaluate this diagnostic on a pinned Qwen2.5 1.5B-to-3B case study
using document-disjoint translator, calibration, and test partitions,
numerical cache controls, online-confidence and target-token-matched baselines,
risk--coverage metrics, and complete batch-one H100 policy timing.  In the
confirmatory run, direct transfer changes the native target's next-token
decision on \result{DIRECT-HARM-PREVALENCE} of test articles; the primary
commutator obtains average precision \result{PRIMARY-COMMUTATOR-AP} versus
\result{PRIMARY-NEGATIVE-MARGIN-AP} for transferred negative margin, while the calibrated
policy reaches coverage \result{PRIMARY-COVERAGE}, risk
\result{PRIMARY-RISK}, and one-sided 95\% Clopper--Pearson upper bound
\result{PRIMARY-RISK-UPPER}.  Its complete expected latency is
\result{PRIMARY-POLICY-LATENCY-MS} ms versus
\result{NATIVE-PREFILL-LATENCY-MS} ms for native prefill.  The predeclared
decision is \result{PRIMARY-DECISION}.  These results concern greedy
next-token compatibility on one model direction; they are not a claim of
agent-action correctness or a distribution-free safety guarantee.
\end{abstract}

\section{Introduction}

Long-running language-model applications increasingly route a conversation
between models of different sizes or specializations.  At a switch, the
receiving model ordinarily prefills the complete accumulated history because
the source model's KV cache lives in a model-specific representation space.
Recent work shows that learned or closed-form mappings can translate caches
across models and, for some pairs, make translation substantially cheaper than
native prefill \citep{fu2026c2c,heo2026crossmodel,lee2026mot}.  The same work
also exposes large variation across model pairs, tasks, and mapper families.
Average benchmark retention therefore does not answer the online systems
question: \emph{is this individual handoff safe to reuse?}

The unavailable counterfactual makes this question difficult.  At handoff
time, a system has the source cache and can translate it.  It does not have the
target model's native cache for the same history; constructing that cache is
the full prefill the transfer is intended to avoid.  Reconstruction error,
attention-output similarity to native state, and native-target logit distance
are useful offline diagnostics but cannot be used as an admission signal on
the serving path.

We investigate whether a short recent suffix can act as a target-native probe.
There are two orders in which translation and suffix processing can be
composed.  First translate the cache of the old prefix and allow the target to
replay the suffix.  Alternatively, allow the source to process the suffix and
then translate the resulting cache.  If the translator preserves the local
state dynamics, these routes should agree.  We call their discrepancy the
\emph{local replay commutator} and ask whether it ranks handoffs for which
direct transfer changes the target's next-token decision.

The distinction between ranking and certification is central.  The
commutator observes a failure of equivariance, not distance to the unknown
native target state.  We give a constructive counterexample in which the
translator commutes exactly with every probed transition while reversing the
target's action.  Thus a positive empirical correlation can support selective
prediction on a frozen distribution, but no result in this paper turns local
commutation into a universal certificate.

Our contributions are:

\begin{enumerate}[leftmargin=*,itemsep=2pt]
  \item We formulate cross-model KV handoff as a per-instance selective
  prediction problem with explicit acceptance coverage, conditional behavior
  risk, fallback cost, and an offline native-target counterfactual.
  \item We define a deployable translation--transition commutator using
  bounded target-side replay, and separate detection of unsafe direct
  transfer from the fidelity of replay as a repair action.
  \item We derive what the commutator measures and prove that even exact
  commutation is insufficient for behavioral fidelity.
  \item We provide a control-heavy experiment that separates online from
  oracle-only quantities, uses one article per example and hash-disjoint
  partitions, and measures complete policy latency rather than component-only
  timing.
  \item We report an evidence-calibrated case study whose empirical outcome is
  \result{PRIMARY-DECISION}; the affine translator is a test instrument, not a
  claimed translation contribution.
\end{enumerate}

\section{Problem formulation}

\subsection{State, transitions, and observable routes}

Let $S$ and $T$ denote source and target autoregressive language models with a
shared tokenizer.  For a token sequence $x$, let $C_m(x)$ be model $m$'s KV
cache after processing $x$.  Let $F_m^r(C)$ append a suffix $r$ to an existing
cache using model $m$, and let $M$ translate a source cache to the target cache
geometry.  Split one fixed-length example into an old prefix $p$, a recent
suffix $r$, and a final anchor token $q$.

The native target logits, used only as an offline oracle, are
\begin{equation}
  \ell_N = G_T(C_T(pr),q),
  \label{eq:native}
\end{equation}
where $G_T(C,q)$ is the target next-token logit vector after consuming anchor
$q$ from cache $C$.  Direct translated-cache inference is
\begin{equation}
  \ell_D = G_T(M(C_S(pr)),q).
  \label{eq:direct}
\end{equation}
The target-replay route translates only the old prefix and processes $r$ with
the target:
\begin{equation}
  \ell_R = G_T(F_T^r(M(C_S(p))),q).
  \label{eq:replay}
\end{equation}
Both $\ell_D$ and $\ell_R$ are online-observable.  Equation~\ref{eq:native} is
not.

The primary harm label is direct-transfer top-one disagreement,
\begin{equation}
  Y_D = \mathbb{1}\!\left[\arg\max \ell_D \ne \arg\max \ell_N\right].
  \label{eq:direct-label}
\end{equation}
We separately define replay disagreement
\begin{equation}
  Y_R = \mathbb{1}\!\left[\arg\max \ell_R \ne \arg\max \ell_N\right].
  \label{eq:replay-label}
\end{equation}
This separation prevents a repaired replay action from being credited with
the fidelity of the direct transfer it is meant to diagnose.

\subsection{Selective reuse}

An online score $s(x)$ induces acceptance rule
$A_\tau(x)=\mathbb{1}[s(x)\leq\tau]$.  Lower scores represent greater
confidence in transfer.  Its acceptance coverage and selective risk are
\begin{align}
  \coverage(\tau) &= \Pr(A_\tau=1), \\
  \risk_D(\tau) &= \Pr(Y_D=1\mid A_\tau=1).
\end{align}
On rejection, the system performs native target prefill.  The intended
operating point is at least 50\% coverage, at most 5\% direct-transfer
disagreement, and at least 30\% measured latency reduction relative to always
prefilling.  These are falsifiable engineering targets, not consequences of
the method definition.

\section{What the commutator can observe}

\subsection{Local translation--transition discrepancy}

At the cache level, the two online routes differ by
\begin{equation}
  \mathcal{K}_r(C_S(p)) =
  F_T^r(M(C_S(p))) - M(F_S^r(C_S(p))).
  \label{eq:commutator}
\end{equation}
The implementation records normalized mean-squared error over the cache
positions written by $r$ and $q$.  Its primary behavioral score is the
Jensen--Shannon divergence between the categorical distributions induced by
$\ell_R$ and $\ell_D$:
\begin{equation}
  s_{\mathrm{JS}}(x)=
  \js\!\left(\operatorname{softmax}(\ell_R),
             \operatorname{softmax}(\ell_D)\right).
  \label{eq:js-score}
\end{equation}

If $M$ were an exact state translation over the probed trajectory, then
$F_T^r\circ M=M\circ F_S^r$ and Equation~\ref{eq:commutator} would vanish.
The converse is false.

\subsection{Error-transport interpretation}

Define the unknown translation error
\begin{equation}
  e(x)=M(C_S(x))-C_T(x).
\end{equation}
Because native target state obeys $F_T^r(C_T(p))=C_T(pr)$, a first-order
expansion of the replay route around $C_T(p)$ gives
\begin{align}
F_T^r(M(C_S(p)))
  &= F_T^r(C_T(p)+e(p)) \\
  &\approx C_T(pr)+J_{F_T^r}e(p).
\end{align}
Meanwhile, direct translation is $M(C_S(pr))=C_T(pr)+e(pr)$.  Therefore
\begin{equation}
  \mathcal{K}_r(C_S(p))
  \approx J_{F_T^r}e(p)-e(pr).
  \label{eq:linearized}
\end{equation}
Direct behavioral harm, however, is governed locally by the target readout,
\begin{equation}
  \ell_D-\ell_N \approx J_{G_T}e(pr).
  \label{eq:behavior-error}
\end{equation}
Equations~\ref{eq:linearized} and~\ref{eq:behavior-error} need not be aligned.
The commutator detects an error-transport residual.  It can miss an error that
is propagated consistently, and it can react strongly to an error direction
that the target output ignores.

\subsection{Exact commutation is not certification}

\paragraph{Proposition 1.}
There exist source state, target state, translator, replay transition, and
target decision functions for which the translation--transition commutator is
zero for every replay probe but direct translated-state behavior differs from
native target behavior.

\paragraph{Proof.}
Let source and target states lie in $\mathbb{R}^2$ and coincide natively:
$C_S(p)=C_T(p)=z=(0,1)$.  Let every source and target replay transition be the
identity, $F_S^r=F_T^r=I$, and let
$M=\operatorname{diag}(1,-1)$.  Then
\[
  F_T^r(Mz)=Mz=M(F_S^r z)
\]
for every $r$, so the commutator is exactly zero.  Let the target decision be
$a(z)=\mathbb{1}[z_2>0]$.  The native decision is $a(z)=1$, while the
translated-state decision is $a(Mz)=0$.  Thus exact commutation does not imply
behavioral fidelity. \hfill$\square$

The proposition rules out a distribution-free certification claim based on
commutation alone.  It does not rule out predictive value on a specified
model pair and input distribution, which is the empirical question tested
below.

\section{Method}

\subsection{Affine cache translators}

The study intentionally uses a small, closed-form translator so that the
verification mechanism remains inspectable.  Qwen2.5 1.5B Instruct is the
source and Qwen2.5 3B Instruct is the target \citep{qwen2024qwen25}.  For
target layer $j$, the depth-aligned source anchor is
\begin{equation}
  a(j)=\operatorname{round}\!\left(
  j\frac{L_S-1}{L_T-1}\right).
  \label{eq:layer-map}
\end{equation}
The primary translator uses a contiguous group of four source layers for each
target layer.  Its start index is
\begin{equation}
  b(j)=\min\!\left(\max(a(j)-1,0),L_S-4\right),
\end{equation}
and its group is $g_4(j)=\{b(j),\ldots,b(j)+3\}$.  Away from the boundaries,
this assigns one layer shallower than the anchor and two deeper, resolving the
even-width tie toward deeper layers.  All KV heads from the four layers are
flattened and concatenated into one feature vector.  A separate affine ridge
map is fit for target keys and values at every target layer with
$\lambda=0.01$.  Sufficient statistics are accumulated in FP64; solved maps
are stored in FP32 and injected into the BF16 target cache.

Rotary position encoding is removed from source and target keys before fitting
and target RoPE is reapplied after mapping \citep{su2024roformer}.  Values are
mapped directly.  This resembles the content-space factorization of
\citet{heo2026crossmodel}, but it is materially weaker than that paper's
production translator: the four layers are fixed by depth and contiguity
rather than selected for predictive utility.  We consequently treat it as a
low-capacity test instrument, not as a representative state-of-the-art
transfer method.  A paired secondary ablation repeats the complete
confirmatory protocol with $g_1(j)=\{a(j)\}$.  The primary claim is always the
contiguous $k=4$ condition; $k=1$ measures dependence on translator capacity.

\subsection{Online scores and actions}

For each replay budget $|r|\in\{8,32,64\}$, we record six scores:

\begin{enumerate}[leftmargin=*,itemsep=2pt]
  \item \textbf{Commutator JS}: Equation~\ref{eq:js-score}, the primary
  behavioral score.
  \item \textbf{Commutator cache NMSE}: normalized cache discrepancy between
  direct and replay routes over the replayed suffix and anchor positions.
  \item \textbf{Transferred negative margin}: the negative difference between
  the largest and second-largest direct-transfer logits.  This is the frozen
  primary baseline.
  \item \textbf{Transferred negative maximum probability}: negative maximum
  softmax probability of the direct-transfer distribution.
  \item \textbf{Transferred entropy}: entropy of the direct-transfer next-token
  distribution.
  \item \textbf{Random}: a deterministic hash of article identity and replay
  length, used as a random-ranking control.
\end{enumerate}

For the direct-transfer label $Y_D$, an accepted policy returns the already
computed direct-transfer action.  For the replay label $Y_R$, it returns the
replay action.  A rejected policy performs fresh native target prefill.  We
also evaluate unconditional replay; if it is both more faithful and cheaper
than selective routing, the proposed policy is dominated.  Replay length 65
is reserved for a target-token-matched control.  At the primary $R=64$, the
commutator computes one direct anchor token plus 65 target tokens on the probe
route, for 66 target-token forwards.  Unconditional replay at $R=65$ also
computes 66 target tokens.  It does not match translation calls, memory work,
or wall time, so the comparison is only target-token-matched.

\section{Experimental protocol}

\subsection{Data and separation}

We use WikiText-103 raw text \citep{merity2017pointer}.  Contiguous dataset rows
are grouped into articles using title markers.  Each article contributes at
most its first 256-token window, so no article contributes multiple correlated
examples.  An SHA-256 digest of article text is reduced modulo three to assign
translator, calibration, and test folds.  The confirmatory run uses the train
split, excludes source rows below 1000, and contains 32 translator articles,
128 calibration articles, and 256 test articles.  Article hashes, raw source
rows, fold identifiers, and token hashes are persisted and audited for
overlap.

Before the confirmatory run, a separate integration preflight uses the
validation split, 128-token windows, and 16/8/16 translator/calibration/test
articles with replay lengths 8 and 32.  Preflight results are not pooled with
the confirmatory estimate.  All sampling is deterministic under seed 1729.
Model and dataset branch names in configuration are resolved to immutable
revisions and stored in the run manifest.

\subsection{Evaluation routes and labels}

Each 256-token example contains 255 context tokens and one anchor token.  The
source processes the 255-token context.  The direct route translates that
cache and lets the target consume only the anchor.  For replay length $R$, the
source cache is sliced after $255-R$ tokens, that prefix is translated, and the
target consumes the last $R$ context tokens plus the anchor.  Native target
prefill over all 256 tokens supplies offline labels and is inaccessible to
online scoring code.

The primary estimand is top-one disagreement in
Equation~\ref{eq:direct-label}.  We additionally retain logit JS divergence,
native top-one margin, replay disagreement, and translated-to-native cache
NMSE as oracle-only diagnostics.  Top-one equality is intentionally a strict
trajectory-fidelity measure: it is neither semantic correctness nor task
success.

\subsection{Calibration and statistics}

For each score and replay length, the threshold is the calibration-score
median.  Calibration labels are not used; the threshold targets 50\% coverage
without adapting to observed errors.  On the untouched test fold we report:

\begin{itemize}[leftmargin=*,itemsep=2pt]
  \item harm prevalence and harm average precision;
  \item risk at 25\%, 50\%, and 75\% descriptive test coverage;
  \item area under the empirical risk--coverage curve;
  \item achieved coverage, conditional risk, and a one-sided exact 95\%
  Clopper--Pearson upper bound at the calibration threshold
  \citep{clopper1934fiducial};
  \item paired bootstrap differences between commutator JS and transferred
  negative margin, using 2000 repetitions and seed 2718.
\end{itemize}

The Clopper--Pearson quantity is an exact binomial test-set summary, not
conformal or distribution-free risk control for future handoffs.  If a strict
deployment guarantee is claimed later, the threshold and risk test should be
replaced with an explicit risk-controlling calibration procedure
\citep{angelopoulos2024conformal}.  The implementation evaluates three replay
budgets and several scores.  The frozen primary condition is commutator JS at
$R=64$ for the full-transfer top-one label, compared with transferred negative
margin by paired bootstrap.  All other score, label, and replay combinations
are secondary or exploratory.

\subsection{Hard controls}

The source and target models independently run the following controls at
sequence lengths 128, 256, and 512 and replay lengths 1, 8, 32, 64, and 65:

\begin{enumerate}[leftmargin=*,itemsep=2pt]
  \item native full prefill versus native prefix cache plus suffix replay;
  \item native cache extraction and reconstruction plus replay;
  \item key de-RoPE/re-RoPE round trip;
  \item same-model identity-cache translation and injection;
  \item SDPA bulk execution versus eager-attention correctness execution.
\end{enumerate}

Native chunking must remain at or below a cache-NMSE limit of $10^{-3}$.
Native full-prefill versus native chunked top-one differences are reported as
the action-noise floor of the two valid execution paths; they do not by
themselves invalidate a run.  Cache reconstruction must be exact.  RoPE and
identity translation are gated against a BF16-derived numerical bound.
Additionally, native full prefill is compared with native chunked replay on
every calibration and test article for every replay length.  Exceeding the
cache-NMSE gate or failing another numerical/evidence control makes the run
invalid and stops performance timing.

\subsection{Timing and evidence}

All timings are batch one on one NVIDIA H100, after explicit warm-up and CUDA
synchronization.  Both wall-clock and CUDA-event durations are stored.  The
confirmatory configuration uses four timing examples per accepted/rejected
stratum, two warm-up executions, and five measured repetitions.  Translation,
cache slicing, online scoring, accepted-action work, and fallback target
prefill are included.  Complete accepted and rejected branches are measured
separately, then weighted by achieved test coverage.  Component timings are
reported only as diagnostics.

The source cache and both models are assumed resident at handoff.  Cache
transport across machines, queueing, batching interactions, and model-loading
latency are outside the measurement boundary.  Each run stores raw JSONL
records before summaries, source/configuration snapshots, environment and
revision metadata, SHA-256 hashes, and an independently recomputed evidence
audit.  Control artifacts are cryptographically bound into the pilot manifest.

\section{Predeclared hypotheses and decisions}

\paragraph{H1: ranking.}
Commutator JS should have higher harm average precision and lower
risk--coverage area than transferred negative margin.  The primary paired
bootstrap interval for AP improvement over negative margin must exclude zero.
Entropy, negative maximum probability, cache NMSE, and random ranking are
secondary comparisons.

\paragraph{H2: operating point.}
The calibration-median policy should attain test coverage of at least 0.5 and
a one-sided 95\% Clopper--Pearson risk upper bound no greater than 0.05.

\paragraph{H3: probe budget.}
Increasing replay from 8 to 32 to 64 tokens should expose more harmful
dynamical mismatch.  Flat or worsening ranking falsifies a simple
more-target-compute explanation.

\paragraph{H4: systems utility.}
The complete coverage-weighted policy should reduce latency by at least 30\%
relative to native target prefill.  Token-equivalent accounting alone cannot
satisfy this hypothesis.

\paragraph{H5: non-domination.}
The selective policy should improve the fidelity--latency frontier over
unconditional replay at the same $R=64$ and over the $R=65$
target-token-matched control.

Runs are classified as \emph{invalid} if a control, partition, or evidence
audit fails; \emph{accept} if the mechanism beats the strongest matched
baseline and reproduces; \emph{revise} if a valid failure isolates one next
intervention; and \emph{reject} if the mechanism is falsified or dominated.

\section{Results}

This section is deliberately machine-fillable.  Bracketed magenta tokens must
be replaced from audited numerical artifacts; they are not estimates.

\subsection{Validity and reproduction}

\begin{table}[t]
\centering
\caption{Artifact identity and validity gates.}
\label{tab:validity}
\resizebox{\columnwidth}{!}{%
\begin{tabular}{lll}
\toprule
Item & Primary run & Reproduction \\
\midrule
Artifact ID & \result{PRIMARY-RUN-ID} & \result{REPRO-RUN-ID} \\
Source control & \result{SOURCE-CONTROL-ID} & \result{SOURCE-CONTROL-REPRO-ID} \\
Target control & \result{TARGET-CONTROL-ID} & \result{TARGET-CONTROL-REPRO-ID} \\
Evidence audit & \result{PRIMARY-AUDIT-STATUS} & \result{REPRO-AUDIT-STATUS} \\
Native chunk top-one floor & \result{PRIMARY-NATIVE-CHUNK-TOP1-FLOOR} & \result{REPRO-NATIVE-CHUNK-TOP1-FLOOR} \\
Maximum native chunk NMSE & \result{PRIMARY-NATIVE-CHUNK-NMSE} & \result{REPRO-NATIVE-CHUNK-NMSE} \\
Record digest match & --- & \result{REPRO-RECORD-MATCH} \\
\bottomrule
\end{tabular}}
\end{table}

The source control's maximum native chunk JS is
\result{SOURCE-MAX-NATIVE-JS}, and maximum cache NMSE is
\result{SOURCE-MAX-NATIVE-CACHE-NMSE}.  The target values are
\result{TARGET-MAX-NATIVE-JS} and
\result{TARGET-MAX-NATIVE-CACHE-NMSE}.  Across confirmatory calibration and
test routes, native chunking changes top one on
\result{CONFIRMATORY-NATIVE-CHUNK-TOP1-FLOOR} of comparisons; this is the
reported action-noise floor.  Its maximum cache NMSE is
\result{CONFIRMATORY-MAX-NATIVE-CHUNK-NMSE}, which is
\result{CONFIRMATORY-NMSE-GATE-STATUS} the $10^{-3}$ hard gate.  The
confirmatory run is therefore \result{RUN-VALIDITY-DECISION}.  If this token
is not ``valid,'' all remaining tables are diagnostic only.

\subsection{Natural transfer harm and ranking}

Direct transfer disagrees with native target top-one on
\result{DIRECT-HARM-COUNT} of \result{TEST-ARTICLE-COUNT} test articles
(\result{DIRECT-HARM-PREVALENCE}).  Table~\ref{tab:ranking} reports ranking
metrics by replay budget.  The frozen primary comparison is commutator JS
versus transferred negative margin at $R=64$ for the full-transfer label.

\begin{table}[t]
\centering
\caption{Test ranking of direct-transfer harm.  AP is harm average precision;
lower AURC is better.}
\label{tab:ranking}
\resizebox{\columnwidth}{!}{%
\begin{tabular}{rrrrrrr}
\toprule
$R$ & Comm. AP & Neg.-margin AP & Entropy AP & Random AP & Comm. AURC & Margin AURC \\
\midrule
8  & \result{R8-COMM-AP}  & \result{R8-MARGIN-AP}  & \result{R8-ENTROPY-AP}  & \result{R8-RANDOM-AP}  & \result{R8-COMM-AURC}  & \result{R8-MARGIN-AURC} \\
32 & \result{R32-COMM-AP} & \result{R32-MARGIN-AP} & \result{R32-ENTROPY-AP} & \result{R32-RANDOM-AP} & \result{R32-COMM-AURC} & \result{R32-MARGIN-AURC} \\
64 & \result{R64-COMM-AP} & \result{R64-MARGIN-AP} & \result{R64-ENTROPY-AP} & \result{R64-RANDOM-AP} & \result{R64-COMM-AURC} & \result{R64-MARGIN-AURC} \\
\bottomrule
\end{tabular}}
\end{table}

At the primary condition, the paired AP difference
$(\mathrm{commutator}-\mathrm{negative\ margin})$ is
\result{PRIMARY-AP-DELTA-MEAN}, with 95\% bootstrap interval
[\result{PRIMARY-AP-DELTA-LOW}, \result{PRIMARY-AP-DELTA-HIGH}].  H1 is
\result{H1-DECISION}.  The replay-budget trend is
\result{REPLAY-BUDGET-TREND}; H3 is \result{H3-DECISION}.

\subsection{Frozen-threshold selective risk}

\begin{table}[t]
\centering
\caption{Policies at calibration-median thresholds.  Upper is the one-sided
exact 95\% Clopper--Pearson bound on conditional top-one disagreement.}
\label{tab:policy}
\resizebox{\columnwidth}{!}{%
\begin{tabular}{rrrrrr}
\toprule
$R$ & Score & Accepted & Coverage & Risk & Upper \\
\midrule
8  & Comm. JS & \result{R8-COMM-ACCEPTED}  & \result{R8-COMM-COVERAGE}  & \result{R8-COMM-RISK}  & \result{R8-COMM-UPPER} \\
8  & Neg. margin & \result{R8-MARGIN-ACCEPTED}   & \result{R8-MARGIN-COVERAGE}   & \result{R8-MARGIN-RISK}   & \result{R8-MARGIN-UPPER} \\
32 & Comm. JS & \result{R32-COMM-ACCEPTED} & \result{R32-COMM-COVERAGE} & \result{R32-COMM-RISK} & \result{R32-COMM-UPPER} \\
32 & Neg. margin & \result{R32-MARGIN-ACCEPTED}  & \result{R32-MARGIN-COVERAGE}  & \result{R32-MARGIN-RISK}  & \result{R32-MARGIN-UPPER} \\
64 & Comm. JS & \result{R64-COMM-ACCEPTED} & \result{R64-COMM-COVERAGE} & \result{R64-COMM-RISK} & \result{R64-COMM-UPPER} \\
64 & Neg. margin & \result{R64-MARGIN-ACCEPTED}  & \result{R64-MARGIN-COVERAGE}  & \result{R64-MARGIN-RISK}  & \result{R64-MARGIN-UPPER} \\
\bottomrule
\end{tabular}}
\end{table}

At the primary condition, the median-calibrated policy accepts
\result{PRIMARY-ACCEPTED-COUNT} examples.  It achieves coverage
\result{PRIMARY-COVERAGE}, risk \result{PRIMARY-RISK}, and upper bound
\result{PRIMARY-RISK-UPPER}.  H2 is \result{H2-DECISION}.  The descriptive
maximum test coverage whose one-sided bound is at most 5\% is
\result{DESCRIPTIVE-SAFE-COVERAGE}; because it uses test labels to select an
operating point, it is not a deployable result.

\subsection{Replay repair and complete cost}

\begin{table}[t]
\centering
\caption{Fidelity and complete batch-one timing.  Policy time is weighted by
full-test accepted coverage using separately measured accepted/rejected
branches.}
\label{tab:timing}
\resizebox{\columnwidth}{!}{%
\begin{tabular}{lrrr}
\toprule
Path & Risk & Wall ms & Savings vs. native \\
\midrule
Native target prefill & 0 & \result{NATIVE-PREFILL-LATENCY-MS} & 0 \\
Direct full transfer & \result{DIRECT-HARM-PREVALENCE} & \result{DIRECT-LATENCY-MS} & \result{DIRECT-LATENCY-SAVINGS} \\
Replay, $R=8$ & \result{R8-REPLAY-RISK} & \result{R8-REPLAY-LATENCY-MS} & \result{R8-REPLAY-LATENCY-SAVINGS} \\
Replay, $R=32$ & \result{R32-REPLAY-RISK} & \result{R32-REPLAY-LATENCY-MS} & \result{R32-REPLAY-LATENCY-SAVINGS} \\
Replay, $R=64$ & \result{R64-REPLAY-RISK} & \result{R64-REPLAY-LATENCY-MS} & \result{R64-REPLAY-LATENCY-SAVINGS} \\
Replay, $R=65$ (target-token-matched) & \result{R65-REPLAY-RISK} & \result{R65-REPLAY-LATENCY-MS} & \result{R65-REPLAY-LATENCY-SAVINGS} \\
Primary selective policy & \result{PRIMARY-RISK} & \result{PRIMARY-POLICY-LATENCY-MS} & \result{PRIMARY-POLICY-LATENCY-SAVINGS} \\
\bottomrule
\end{tabular}}
\end{table}

Unconditional replay changes the target top-one on
\result{PRIMARY-REPLAY-RISK} of examples at the primary replay length.  The
selective policy is \result{SAME-R64-DOMINATION-STATUS} relative to this
same-length repair baseline and \result{R65-DOMINATION-STATUS} relative to the
$R=65$ target-token-matched control.  The latter matches 66 target-token
forwards but not total work.  H4 is \result{H4-DECISION}, and H5 is
\result{H5-DECISION}.  CUDA-event results lead to
\result{CUDA-TIMING-CONSISTENCY} the wall-clock conclusion.

\subsection{Error structure}

Among direct-transfer failures, the primary score's false-accept set contains
\result{FALSE-ACCEPT-COUNT} articles.  Their median native margin is
\result{FALSE-ACCEPT-NATIVE-MARGIN}, versus
\result{TRUE-REJECT-NATIVE-MARGIN} for correctly rejected harmful transfers.
Replay repairs \result{REPLAY-REPAIR-COUNT} direct failures and introduces
\result{REPLAY-INTRODUCED-ERROR-COUNT} errors among otherwise faithful direct
transfers.  The observed pattern is \result{ERROR-STRUCTURE-SUMMARY}.  This
analysis is descriptive and does not alter the frozen score or threshold.

\subsection{Paired translator-capacity ablation}

The secondary $k=1$ run uses the same resolved models, article identities,
partitions, replay budgets, seed, ridge penalty, scores, labels, and timing
procedure as the primary $k=4$ run.  Only the number of contiguous source
layers supplied to each target-layer regression changes.  Table~\ref{tab:k-ablation}
therefore tests whether the primary conclusion is specific to translator
capacity; it is not a second confirmatory hypothesis.

\begin{table}[t]
\centering
\caption{Paired translator-capacity ablation at the frozen $R=64$ condition.
Upper is the one-sided 95\% Clopper--Pearson bound.}
\label{tab:k-ablation}
\resizebox{\columnwidth}{!}{%
\begin{tabular}{rrrrrr}
\toprule
Source layers & Transfer harm & Comm. AP & Margin AP & Coverage & Risk / Upper \\
\midrule
$k=4$ & \result{K4-TRANSFER-HARM} & \result{K4-COMM-AP} & \result{K4-MARGIN-AP} & \result{K4-COVERAGE} & \result{K4-RISK}~/~\result{K4-CP-UPPER} \\
$k=1$ & \result{K1-TRANSFER-HARM} & \result{K1-COMM-AP} & \result{K1-MARGIN-AP} & \result{K1-COVERAGE} & \result{K1-RISK}~/~\result{K1-CP-UPPER} \\
\bottomrule
\end{tabular}}
\end{table}

Moving from $k=4$ to $k=1$ changes commutator AP by
\result{K1-MINUS-K4-COMM-AP}, transfer-harm prevalence by
\result{K1-MINUS-K4-HARM}, and complete policy latency by
\result{K1-MINUS-K4-POLICY-LATENCY-MS} ms.  The capacity dependence is
\result{TRANSLATOR-CAPACITY-INTERPRETATION}.

\subsection{Decision}

The confirmatory trial decision is \textbf{\result{PRIMARY-DECISION}} because
\result{DECISION-RATIONALE}.  A fresh-process reproduction is
\result{REPRODUCTION-DECISION}.  Accordingly, the evidence
\result{SUPPORTS-OR-DOES-NOT-SUPPORT} the claim that one bounded local replay
probe is sufficient for the target operating point.  Regardless of this
empirical decision, Proposition~1 establishes that commutation alone cannot be
a universal certificate.

\section{Related work}

\paragraph{Cross-model cache translation.}
Cache-to-Cache learns neural projectors and fusers for latent semantic
communication between heterogeneous models \citep{fu2026c2c}.  Latent Cache
Flow compresses translated cache summaries and permits sender and receiver to
hold different contexts \citep{rossi2026lcf}.  These methods optimize
communication utility rather than reproducing the receiver's native prefill.
\citet{heo2026crossmodel} directly targets prefill reuse with closed-form
per-head ridge maps, cross-layer source selection, and RoPE-stripped key
mapping.  It shows substantial pair-level variation and that attention-aligned
error placement predicts downstream retention better than global
reconstruction fit.  Those diagnostics compare with native target quantities
or summarize a model pair; our question is per-handoff admission without the
native target cache.  Mixture-of-Translators uses token-routed translators and
a Context Correction Loss that aligns a replayed target trajectory with the
native target trajectory during training \citep{lee2026mot}.  Our use of
replay is different: the translator remains fixed, native target state is
unavailable online, and path disagreement is evaluated as an admission score.

\paragraph{Partial recomputation and cache reuse.}
DroidSpeak shares KV state across fine-tuned variants with identical
architecture and selectively recomputes critical layers
\citep{liu2024droidspeak}.  CacheBlend selectively recomputes tokens to restore
cross-chunk interactions when combining same-model RAG caches
\citep{yao2024cacheblend}; ProphetKV makes that selection query-dependent
\citep{wang2026prophetkv}.  These systems establish recomputation as a useful
repair budget, but their cached states do not arise from an independently
trained cross-model translator.  We therefore include unconditional recent
replay as the strongest direct systems baseline rather than treating replay
cost as free.

\paragraph{Behavioral consistency under reuse.}
\citet{liang2026reusefails} show that reuse can preserve aggregate task
accuracy while substantially changing an LLM judge's selected candidate.
Their result motivates comparison to dense-prefill decisions rather than task
averages alone and prevents a broad claim that this is the first risk-aware
study of cache reuse.  Our narrower contribution concerns online prediction
of per-instance error from cross-model translated state.

\paragraph{Selective prediction and risk control.}
Selective classification formalizes the trade-off between conditional error
and accepted coverage \citep{geifman2019selectivenet}.  Conformal risk control
provides finite-sample control for suitable calibrated losses
\citep{angelopoulos2024conformal}.  We borrow risk--coverage terminology but do
not claim conformal validity: the current threshold is an unlabeled median and
the reported test interval is an exact one-sided binomial interval.

\section{Discussion}

\subsection{How to interpret a positive ranking result}

If H1 succeeds while H2 fails, the commutator contains information but is not
selective enough at the desired operating point.  That result motivates richer
probe sets or a learned router; it does not justify deployment.  If H2 succeeds
but H4 fails at 256 tokens, the verifier may be behaviorally useful yet
economically misplaced.  Translation and replay have fixed overheads, whereas
target prefill becomes more expensive with context length.  The next systems
experiment should therefore freeze the fidelity policy and measure its
latency crossover over longer contexts, not retune the score at each length.

\subsection{How to interpret a negative result}

A failed commutator score has at least three possible causes.  First, harmful
translation error may be a shared mode visible to neither route, as in
Proposition~1.  Second, the replay suffix may not excite an error direction
relevant to the anchor decision.  Third, even the primary contiguous
four-layer affine translator may place the experiment in a regime where almost
every state is damaged and ranking is statistically or operationally
unhelpful.  The paired $k=1$ ablation can reveal capacity sensitivity, but it
cannot establish behavior under a stronger learned translator.  Distinguishing
the remaining explanations requires stronger translator-quality regimes and
multiple, frozen probes.  It cannot be resolved by selecting the best
test-set threshold after the fact.

\subsection{Required next evidence}

A claim about agent handoffs requires a structured action dataset.  The native
and transferred outputs should be decoded deterministically into a canonical
pair of tool name and arguments, with failures grouped by full trajectory so
that related calls cannot cross partitions.  A claim about general cross-model
reuse additionally requires several source--target directions, at least one
strong multi-source or nonlinear translator, and a long-context timing sweep.
These are extensions, not claims of the present case study.

\section{Limitations}

\begin{enumerate}[leftmargin=*,itemsep=2pt]
  \item \textbf{One direction.}  The empirical study covers only Qwen2.5
  1.5B Instruct to 3B Instruct.  Models share a tokenizer, family, compatible
  rotary structure, and KV geometry.  Results do not establish cross-family or
  mismatched-KV behavior.
  \item \textbf{Restricted translators.}  The primary translator concatenates
  four fixed contiguous depth-aligned source layers; the secondary ablation
  uses one.  Neither selects source layers for predictive utility nor trains a
  nonlinear mapper.  Failure here cannot establish failure for a
  state-of-the-art translator \citep{heo2026crossmodel}.
  \item \textbf{Text-only, one-step label.}  WikiText next-token top-one
  agreement is a trajectory-fidelity canary, not semantic quality, agent
  success, or tool-action equality.  Equal first tokens can lead to different
  later trajectories, while a near-tie flip can be semantically harmless.
  \item \textbf{Short context.}  The confirmatory context length is 256.
  Fixed translation and probe overhead may dominate in this regime.  No
  long-context speedup is implied without a measured crossover sweep.
  \item \textbf{Distribution dependence.}  Calibration and test articles are
  independent under a deterministic WikiText partition but share a domain.
  No robustness to domain or temporal shift is measured.
  \item \textbf{No formal risk guarantee.}  A calibration median targets
  coverage, not error.  The one-sided Clopper--Pearson bound summarizes an
  untouched test sample; it is not a distribution-free admission guarantee.
  \item \textbf{Systems boundary.}  Timing assumes resident source cache and
  resident models on one H100 and excludes inter-node transfer, serving-engine
  scheduling, batching, and cache serialization.
  \item \textbf{Multiplicity.}  Three replay budgets and multiple scores are
  computed.  Only a condition fixed before confirmatory test inspection can
  support a primary inferential claim.
\end{enumerate}

\section{Conclusion}

Cross-model KV translation creates a selective-computation problem: the target
should reuse a translated state only when the saved prefill is worth its
behavioral risk.  Local replay supplies a natural online self-consistency
signal because a faithful translator should approximately intertwine source
and target transitions.  Yet commutation and correctness are different
properties.  The formal counterexample shows that even perfect local
commutation can coexist with a changed target decision.  The audited case
study yields decision \result{PRIMARY-DECISION}: \result{CONCLUSION-SENTENCE}.
The appropriate conclusion is therefore bounded to instance-level next-token
prediction on the tested pair.  Stronger claims require stronger translators,
structured actions, multiple model directions, explicit risk calibration, and
long-context end-to-end timing.

\bibliographystyle{plainnat}
\bibliography{refs}

\appendix

\section{Exact confirmatory configuration}

\begin{table}[h]
\centering
\caption{Frozen experiment configuration.  Immutable revisions are populated
from the audited manifest.}
\label{tab:config}
\resizebox{\columnwidth}{!}{%
\begin{tabular}{ll}
\toprule
Field & Value \\
\midrule
Source & Qwen/Qwen2.5-1.5B-Instruct @ \result{SOURCE-REVISION-SHA} \\
Target & Qwen/Qwen2.5-3B-Instruct @ \result{TARGET-REVISION-SHA} \\
Attention backend & SDPA; eager correctness control \\
Precision & BF16 model/cache; FP64 ridge statistics; FP32 map \\
Dataset & Salesforce/wikitext, wikitext-103-raw-v1 @ \result{DATASET-REVISION-SHA} \\
Dataset split & train \\
Window length & 256 tokens (255 context + 1 anchor) \\
Translator/calibration/test & 32 / 128 / 256 articles \\
Article folds & SHA-256 modulo 3: 0 / 1 / 2 \\
Minimum raw row & 1000 for every partition \\
Replay lengths & 8, 32, 64; control 65 \\
Primary translator & $k=4$ contiguous anchored groups \\
Secondary translator & paired $k=1$ depth-aligned ablation \\
Primary analysis & commutator JS, $R=64$, full-transfer label \\
Primary comparator & transferred negative margin \\
Target-token-matched control & unconditional replay, $R=65$ \\
Ridge penalty & 0.01 \\
Target coverage / risk target & 0.50 / 0.05 \\
Bootstrap & 2000 repetitions, seed 2718 \\
Sampling seed & 1729 \\
Timing & 4 examples/stratum, 2 warm-up, 5 repeats \\
Hardware & \result{GPU-IDENTITY} \\
Software manifest & \result{ENVIRONMENT-ARTIFACT-SHA} \\
\bottomrule
\end{tabular}}
\end{table}

The integration preflight uses the validation split, length 128, 16/8/16
articles, replay lengths 8 and 32, 200 bootstrap repetitions, two timing
examples, one warm-up, and two measured repeats.  It validates plumbing only
and contributes no observations to Tables~\ref{tab:ranking}--\ref{tab:timing}.

\section{Evidence contract}

A valid paper result must be traceable to the following artifact fields:

\begin{itemize}[leftmargin=*,itemsep=2pt]
  \item resolved model and dataset revisions;
  \item configuration, source snapshot, environment, data, translator, raw
  record, timing-record, and summary SHA-256 digests;
  \item source and target control artifact IDs and their recomputed summaries;
  \item one raw record for every partition/article/replay-length Cartesian
  product expected by configuration;
  \item disjoint document hashes, token hashes, and source rows across
  translator, calibration, and test;
  \item exact recomputation of summary metrics from JSONL records;
  \item complete timing records for native, direct, replay, score overhead,
  accepted policy branches, and rejected policy branches;
  \item a fresh-process reproduction whose deterministic raw-record digest is
  \result{REPRO-RECORD-MATCH} the primary run.
\end{itemize}

No plot or number copied from an interactive notebook is authoritative unless
it can be regenerated from these artifacts.

\section{Reporting boundaries}

The following language is permitted if supported by filled artifacts:
``on the pinned Qwen2.5 direction and WikiText partition, the online score
ranked next-token disagreement with metric X.''  The following language is not
supported: ``the verifier makes cross-model agent handoffs safe,'' ``the method
guarantees at most 5\% errors,'' or ``the system accelerates long-context
serving.''  Native-target fidelity is a counterfactual-consistency target, not
ground-truth answer correctness.

\end{document}
